\chapter{Umbilical hernia}
\index{Umbilical hernia}

\section*{Definition and Overview}
An umbilical hernia represents a protrusion of intra-abdominal contents through a defect in the umbilical ring. The umbilical ring is a natural opening in the abdominal wall that, during fetal development, allowed the passage of the umbilical vessels. In normal physiology, this ring closes shortly after birth, forming the umbilicus, commonly referred to as the belly button. However, when the surrounding fibromuscular tissue fails to seal completely, or when the umbilical scar is weakened later in life, abdominal contents---such as preperitoneal fat, omentum, or loops of the small and large intestines---can herniate through the gap. This results in a visible or palpable bulge in the umbilical region.

Umbilical hernias are classified based on the age of onset into congenital (infantile) hernias and acquired (adult) hernias. While congenital umbilical hernias are extremely common in neonates and often resolve spontaneously as the abdominal wall matures, acquired umbilical hernias in adults are progressive conditions that rarely resolve without surgical intervention. Understanding the distinction between these two forms is essential for clinical management, as the treatment paths diverge significantly. This chapter explores the anatomical basis, clinical characteristics, diagnostic approaches, and therapeutic options for both paediatric and adult umbilical hernias, underscoring the importance of early detection to prevent potentially life-threatening complications.

\section*{Epidemiology}
The epidemiology of umbilical hernias varies markedly between paediatric and adult populations.
In infants, umbilical hernias are highly prevalent, occurring in approximately 10\% to 25\% of all newborns. The incidence is significantly higher in specific subgroups. Premature infants and low-birth-weight babies are particularly susceptible, with prevalence rates reaching up to 75\% to 80\% in infants weighing less than 1,200 grams at birth. Demographically, congenital umbilical hernias are observed with equal frequency in male and female infants. However, a strong racial disparity exists, with children of African descent displaying a five-to-eight-fold higher incidence compared to Caucasian children. The biological basis for this racial variation remains poorly understood, but genetic factors are believed to play a significant role.

In adults, umbilical hernias account for approximately 10\% of all abdominal wall hernias. Unlike the paediatric cohort, acquired umbilical hernias are far more common in females than in males, with a female-to-male ratio of roughly three to one. This disparity is primarily attributed to the anatomical and physiological changes associated with pregnancy. The incidence peaks during middle age, particularly in the fourth, fifth, and sixth decades of life. Obese individuals and multiparous women represent the highest-risk adult demographics, reflecting the causative role of chronic intra-abdominal pressure in the pathogenesis of adult hernias.

\section*{Aetiology and Pathophysiology}
The underlying causes and physiological mechanisms of umbilical hernias differ fundamentally between children and adults. In infants, the primary mechanism is developmental, arising from the incomplete closure of the umbilical ring. During gestation, the umbilical cord vessels pass through a defect in the linea alba. Following birth and the clamping of the umbilical cord, these vessels thrombose and fibrose, forming the umbilical ligaments. The surrounding tissue is expected to contract and close the ring. A failure in this physiological contraction leaves a patent fascial opening.

In adults, the hernia is typically acquired and results from a combination of anatomical weakness and chronic elevation of intra-abdominal pressure. The umbilical area is inherently a site of relative weakness in the linea alba. When continuous or repetitive pressure is exerted against this area, the thin fascial layers stretch and eventually tear, creating a defect. The pathophysiological processes can be categorized into the following major risk factors and mechanisms:
\begin{itemize}
    \item \textbf{Congenital Fascial Defect:} In infants, the failure of the lateral abdominal wall muscles to fuse completely at the midline during the neonatal period leaves the umbilical ring patent.
    \item \textbf{Obesity:} Excess visceral and subcutaneous abdominal fat increases intra-abdominal pressure, stretching the abdominal wall and widening the umbilical ring.
    \item \textbf{Pregnancy:} Multiple pregnancies or carrying twins or triplets significantly distends the abdominal wall, weakens the linea alba, and expands the umbilical defect.
    \item \textbf{Ascites:} The accumulation of fluid within the peritoneal cavity, commonly secondary to liver cirrhosis, heart failure, or renal disease, creates persistent hydrostatic pressure that forces abdominal contents through the umbilical ring.
    \item \textbf{Chronic Coughing or Straining:} Conditions such as chronic obstructive pulmonary disease (COPD), chronic constipation, or benign prostatic hyperplasia (BPH) lead to repetitive, acute increases in abdominal pressure, promoting fascial degradation.
    \item \textbf{Heavy Lifting and Strenuous Labour:} Repetitive lifting of heavy loads causes intense, sudden increases in intra-abdominal pressure, which can tear vulnerable fascial fibres over time.
    \item \textbf{Connective Tissue Disorders:} Genetic conditions that affect collagen synthesis and tissue elasticity, such as Marfan syndrome or Ehlers-Danlos syndrome, compromise the structural integrity of the linea alba and predispose individuals to hernia formation.
\end{itemize}

\section*{Clinical Features}
The clinical presentation of an umbilical hernia is characterized by a focal swelling or bulge at or immediately adjacent to the navel. The symptoms and physical signs vary depending on the size of the defect and whether the hernia is reducible or incarcerated.
\begin{itemize}
    \item \textbf{Visible and Palpable Bulge:} The most common feature is a soft protrusion at the umbilicus. This bulge may range in size from less than a centimetre to several centimetres in diameter.
    \item \textbf{Fluctuation with Pressure Changes:} In both infants and adults, the bulge typically becomes more prominent during activities that increase intra-abdominal pressure, such as coughing, crying, straining, lifting, or standing upright. Conversely, the bulge may decrease in size or disappear completely when the individual lies flat in a supine position.
    \item \textbf{Reducibility:} A reducible hernia can be gently pushed back into the abdominal cavity. In infants, these hernias are almost always reducible and rarely cause pain. When the infant cries, the hernia may swell, but it returns to its normal state when the child is calm.
    \item \textbf{Localized Pain and Discomfort:} In adults, umbilical hernias are frequently accompanied by a dull ache, dragging sensation, or localized tenderness, particularly after prolonged standing or physical exertion.
    \item \textbf{Skin Changes:} The overlying skin of a chronic, large hernia may become stretched, thinned, or discoloured. In severe cases, friction from clothing can lead to ulceration, dermatitis, or skin breakdown over the apex of the hernia.
    \item \textbf{Incarceration and Strangulation Signs:} If the herniated contents become trapped outside the abdominal wall, the bulge becomes firm, tender, and non-reducible. If the blood supply to the trapped tissue is compromised, the bulge may turn purple, red, or dark blue, accompanied by severe, constant pain, nausea, vomiting, and systemic signs of bowel obstruction.
\end{itemize}

\section*{Differential Diagnosis}
It is crucial to distinguish an umbilical hernia from other pathological conditions that present with masses or swelling in the umbilical and periumbilical regions. The differential diagnosis includes:
\begin{itemize}
    \item \textbf{Paraumbilical Hernia:} This is a protrusion through a defect in the linea alba immediately adjacent to (usually above or below) the umbilical ring, rather than directly through the ring itself. Although clinically managed similarly to umbilical hernias in adults, they have a higher risk of incarceration because the fascial defect is often smaller and more rigid.
    \item \textbf{Epigastric Hernia:} A hernia occurring in the midline between the umbilicus and the xiphoid process. It is caused by a defect in the linea alba and is often smaller and more painful than an umbilical hernia.
    \item \textbf{Umbilical Granuloma:} A common cause of umbilical swelling in neonates, presenting as a soft, pinkish-red nodule that forms during the healing process after the umbilical cord detaches. Unlike a hernia, it does not fluctuate with crying or intra-abdominal pressure and often oozes serosanguinous fluid.
    \item \textbf{Patent Urachus or Omphalomesenteric Duct Cyst:} Congenital anomalies where embryological structures fail to obliterate. They may present as an umbilical discharge, mass, or recurrent local infection, requiring distinct surgical management.
    \item \textbf{Sister Mary Joseph Nodule:} A metastatic lesion of the umbilicus, secondary to an advanced intra-abdominal malignancy (typically gastric, ovarian, colorectal, or pancreatic cancer). It presents as a hard, irregular, and sometimes ulcerated nodule, and must be differentiated from a benign hernia.
    \item \textbf{Caput Medusae:} Dilated, tortuous cutaneous veins radiating from the umbilicus, commonly seen in patients with severe portal hypertension (such as in liver cirrhosis). While it causes a visible change in the umbilical area, it is vascular rather than structural and lacks a fascial defect.
    \item \textbf{Primary Umbilical Tumours or Endometriosis:} Rare benign or malignant neoplasms (e.g., lipoma, fibroma, or melanoma) or deposits of endometrial tissue (umbilical endometriosis, also known as Villar's nodule) that present as a painful, firm nodule at the umbilicus, often swelling or bleeding cyclically in the case of endometriosis.
\end{itemize}

\section*{When to Seek Medical Care}
While many umbilical hernias are benign and can be managed conservatively or scheduled for elective surgery, certain symptoms indicate a medical emergency. Delaying care in these situations can lead to irreversible tissue necrosis and life-threatening complications.

Patients or caregivers should contact a primary care physician if they notice a new, persistent bulge at the navel, or if an existing hernia increases in size or becomes consistently uncomfortable.

Immediate emergency medical attention must be sought, or emergency services called (such as \textbf{911} in the United States, \textbf{999} in the United Kingdom, or \textbf{local emergency services} in many countries), if any of the following symptoms develop:
\begin{itemize}
    \item The hernia bulge becomes hard, extremely tender to the touch, and cannot be gently pushed back into the abdomen (non-reducible or incarcerated).
    \item The skin overlying the hernia turns red, purple, dark blue, or black, indicating potential ischemia or tissue death.
    \item The patient experiences sudden, severe, and worsening pain at the hernia site.
    \item The onset of systemic signs of bowel obstruction, including intractable nausea, projectile vomiting, and an inability to pass gas or have a bowel movement.
    \item The development of systemic signs of infection or shock, such as a high fever, rapid heart rate, confusion, or a sudden drop in blood pressure.
\end{itemize}

\section*{Diagnosis and Investigations}
The diagnosis of an umbilical hernia is primarily clinical, relying on a detailed medical history and a comprehensive physical examination. Diagnostic imaging is reserved for specific circumstances.
\begin{itemize}
    \item \textbf{Clinical History:} The physician inquires about the onset, duration, and progression of the swelling, the presence of pain or discomfort, and potential risk factors such as heavy lifting, chronic cough, multiple pregnancies, or a history of ascites.
    \item \textbf{Physical Examination:} The examination is conducted with the patient in both the supine (lying down) and standing positions. The doctor inspects the umbilical region for asymmetry, swelling, and skin changes. Palpation is performed to assess the size of the fascial defect, the consistency of the herniated contents, and whether the hernia is easily reducible. The patient may be asked to cough or perform a Valsalva manoeuvre to make the hernia bulge more prominent.
    \item \textbf{Ultrasonography (US):} Ultrasound is the first-line imaging modality of choice, particularly in pregnant women, children, and patients with significant overlying adipose tissue where physical examination is challenging. It is non-invasive, avoids radiation, and can accurately measure the size of the fascial defect, identify the herniated contents (fat versus bowel), and assess the patency of the blood supply.
    \item \textbf{Computed Tomography (CT) Scan:} A CT scan of the abdomen and pelvis is indicated in complex, large, or recurrent hernias, or when there is suspicion of acute complications such as incarceration, strangulation, or bowel obstruction. It provides detailed anatomical views of the abdominal wall and helps in surgical planning.
    \item \textbf{Magnetic Resonance Imaging (MRI):} MRI is rarely used but may be considered in pregnant patients with complex presentations when ultrasound is inconclusive, as it provides detailed tissue characterization without ionizing radiation.
    \item \textbf{Laboratory Investigations:} Routine blood tests, including a complete blood count (CBC), electrolytes, renal function tests, and inflammatory markers (such as C-reactive protein), are not required for simple hernias but are essential if strangulation, bowel necrosis, or systemic infection is suspected.
\end{itemize}

\section*{Management}
The management strategy for an umbilical hernia depends on the patient's age, the severity of the symptoms, the size of the fascial defect, and the risk of complications.

In paediatric patients, a conservative approach is the standard of care. Because the vast majority of congenital umbilical hernias close spontaneously by the age of three to five years as the abdominal wall muscles strengthen, active intervention is rarely needed. Parents are reassured, and the hernia is monitored during regular wellness visits. Surgical repair in children is indicated only if the hernia is very large (greater than 1.5 to 2 centimetres in diameter), fails to close by age four or five, becomes incarcerated, or causes cosmetic distress to the child.

In adult patients, watch-and-wait management is generally not recommended, except for patients with very small, asymptomatic hernias or those with significant comorbidities that make surgery highly risky. Because adult umbilical hernias tend to enlarge over time and have a higher risk of incarceration (roughly 5\% to 10\%), elective surgical repair is the preferred treatment. Management options include:
\begin{itemize}
    \item \textbf{Conservative Management:} This involves lifestyle modifications and observation. In patients unfit for surgery, an abdominal binder or truss may be used to provide external support and prevent the bulge from protruding, though this does not repair the defect.
    \item \textbf{Primary Suture Repair (Herniorrhaphy):} For small fascial defects (typically less than 1 to 2 centimetres in diameter), the surgical repair involves exposing the hernia sac, reducing the contents, and closing the fascial edges directly using strong, non-absorbable sutures.
    \item \textbf{Mesh Repair (Hernioplasty):} For defects larger than 2 centimetres, primary suture repair carries a high recurrence rate. In these cases, a synthetic or biological mesh is placed over or under the fascial defect and secured. The mesh acts as a scaffold for new tissue growth, distributing tension across the abdominal wall and significantly reducing the risk of recurrence.
    \item \textbf{Open Surgery:} The traditional surgical approach involves making a small incision at the base of the umbilicus, dissecting the hernia sac, and repairing the fascia. It is highly effective and widely used.
    \item \textbf{Laparoscopic and Robotic Surgery:} Minimally invasive techniques involve inserting a camera and specialized instruments through small, distant incisions. These methods are particularly beneficial for very large hernias, recurrent hernias, or obese patients. They are associated with less postoperative pain, lower wound infection rates, and faster recovery times.
    \item \textbf{Management of Ascites:} In patients where the hernia is driven by ascites (e.g., decompensated cirrhosis), control of the fluid accumulation through low-sodium diets, diuretics, or therapeutic paracentesis is mandatory before and after surgery to prevent recurrence and wound healing complications.
\end{itemize}

\section*{Complications}
Untreated umbilical hernias, particularly in adults, can lead to serious and potentially life-threatening complications. Additionally, surgical treatment itself carries inherent risks.
\begin{itemize}
    \item \textbf{Incarceration:} This occurs when the herniated tissue (often omentum or a loop of intestine) becomes trapped within the hernia sac and cannot be pushed back into the abdominal cavity. Incarceration causes localized pain and is the precursor to strangulation.
    \item \textbf{Strangulation:} A critical medical emergency where the blood supply to the incarcerated tissue is cut off due to constriction at the neck of the hernia. This leads to rapid ischemia, tissue necrosis, gangrene, and potential perforation of the bowel, which can cause life-threatening peritonitis and systemic sepsis.
    \item \textbf{Intestinal Obstruction:} If a loop of the small or large bowel is trapped within the hernia, it can block the passage of food, fluids, and gas. This results in severe abdominal pain, bloating, projectile vomiting, and an inability to pass stool.
    \item \textbf{Surgical Site Infections (SSI):} Postoperative complications include superficial wound infections or deep mesh infections. Mesh infections are particularly challenging and may require long-term antibiotic therapy or surgical removal of the mesh.
    \item \textbf{Seroma and Haematoma Formation:} The accumulation of fluid (seroma) or blood (haematoma) in the dead space created during the surgical dissection. While usually self-limiting, large accumulations may require aspiration or surgical drainage.
    \item \textbf{Hernia Recurrence:} Despite successful surgical repair, there is a risk that the hernia will return. Recurrence rates are higher with primary suture repair of large defects, poor surgical technique, wound infections, uncontrolled ascites, persistent obesity, or smoking.
    \item \textbf{Chronic Pain:} Persistent discomfort or neuropathic pain at the surgical site can occur, often secondary to nerve entrapment by sutures or mesh, or chronic inflammatory reactions.
\end{itemize}

\section*{Prognosis and Outcomes}
The overall prognosis for patients with an umbilical hernia is excellent, particularly when the condition is identified early and managed appropriately.

In children, the prognosis is outstanding. More than 90\% of congenital umbilical hernias resolve spontaneously without surgical intervention before the age of five. For those children who do require surgery, the procedure is straightforward, with minimal postoperative complications and a recurrence rate of less than 1\%. Long-term recovery is rapid, and children typically return to normal activities within a few days.

In adults, the prognosis is highly favourable for those undergoing elective, uncomplicated surgical repair. The recurrence rate is low, particularly when mesh is utilized, ranging from 2\% to 5\% compared to up to 10\% to 20\% for suture-only repairs of larger defects. Recovery time varies based on the surgical approach; patients undergoing laparoscopic or robotic repair often return to light activities within one to two weeks, whereas open repairs may require four to six weeks of restricted physical activity.

However, the prognosis worsens significantly if the patient presents with acute complications such as strangulation or bowel necrosis. Emergency hernia repairs carry a much higher morbidity and mortality rate, particularly in elderly patients or those with severe underlying conditions such as liver failure or advanced cardiac disease. Early intervention is therefore the primary factor in ensuring optimal outcomes.

\section*{Prevention and Self-Care}
While congenital umbilical hernias in infants cannot be prevented, there are practical steps that adults can take to reduce the risk of developing an acquired hernia, prevent the worsening of an existing one, and support recovery after surgical repair.
\begin{itemize}
    \item \textbf{Maintain a Healthy Body Weight:} Reducing excess body weight helps lower the chronic, resting pressure exerted on the abdominal wall, decreasing the likelihood of fascial stretching and herniation.
    \item \textbf{Adopt Proper Lifting Techniques:} When lifting heavy objects, individuals should bend at the knees and use the strength of their legs rather than bending at the waist and straining the abdominal muscles. Exhaling during the lift rather than holding one's breath also helps minimize intra-abdominal pressure spikes.
    \item \textbf{Manage Chronic Coughing:} Seeking medical treatment for chronic respiratory conditions, such as asthma or chronic bronchitis, and quitting smoking to prevent smoker's cough, are essential preventive measures.
    \item \textbf{Prevent Constipation:} Eating a high-fibre diet rich in fruits, vegetables, and whole grains, staying adequately hydrated, and staying physically active can prevent the need to strain during bowel movements. Over-the-counter stool softeners may be used if recommended by a physician.
    \item \textbf{Use External Abdominal Support:} For patients with small hernias who are not immediate candidates for surgery, wearing a supportive abdominal binder or truss can help contain the bulge during physical activity, though it must be used under medical supervision.
    \item \textbf{Adhere to Postoperative Restrictions:} After undergoing surgical repair, patients must strictly follow their surgeon's instructions. This includes avoiding heavy lifting (usually defined as anything over 5 kilograms) and avoiding strenuous exercise for four to six weeks to allow the fascial closure and mesh to heal securely.
    \item \textbf{Monitor for Warning Signs:} Patients with a known hernia or those recovering from surgery should routinely inspect the umbilical area for changes in size, colour, or tenderness and report any concerns to their healthcare provider immediately.
\end{itemize}

\section*{Sources and Further Reading}
To obtain further reliable and detailed information regarding umbilical hernias, their diagnosis, and treatment options, the following resources are recommended:
\begin{itemize}
    \item \textbf{National Health Service (NHS):} Umbilical Hernia Repair Overview --- \url{https://www.nhs.uk/conditions/umbilical-hernia-repair/}
    \item \textbf{Mayo Clinic:} Umbilical Hernia Symptoms and Causes --- \url{https://www.mayoclinic.org/diseases-conditions/umbilical-hernia/symptoms-causes/syc-20378685}
    \item \textbf{Cleveland Clinic:} Umbilical Hernia Diagnosis and Management --- \url{https://my.clevelandclinic.org/health/diseases/14381-umbilical-hernia}
    \item \textbf{American College of Surgeons (ACS):} Hernia Repair Patient Education --- \url{https://www.facs.org/for-patients/safe-surgery/hernia/}
    \item \textbf{Johns Hopkins Medicine:} Pediatric Umbilical Hernia Care --- \url{https://www.hopkinsmedicine.org/health/conditions-and-diseases/hernias/umbilical-hernia-in-children}
    \item \textbf{Healthdirect Australia:} Umbilical Hernia Information --- \url{https://www.healthdirect.gov.au/umbilical-hernia}
\end{itemize}